\documentclass[11pt]{article}

\usepackage[margin=1in]{geometry}

\usepackage{amsmath}

\usepackage[T1]{fontenc}

\title{Mathematical Modeling and Numerical Simulation Supplement}

\author{Anonymous Research Plan Author}

\date{}

\begin{document}

\setlength{\emergencystretch}{3em}

\maketitle

\section{Q1 (final: v0)}

\noindent\textbf{Abstract—} A parsimonious ordinary differential equation framework modeling the coupled evolution of neutron star spin period and magnetic inclination angle under vacuum dipole braking, wind enhancement, inclination decay, and accretion-driven spin-up. The model isolates the dynamical signatures of isolated magnetic-dipole evolution versus binary recycling by encoding phase-dependent torque switching directly into the state derivatives, enabling population-level statistical comparison against observed braking indices and characteristic ages.

\subsection{Question and Scope}

\noindent\textbf{Question.} Can the joint distribution of braking index n and characteristic age tau\_c statistically distinguish isolated magnetic-dipole evolution from binary recycling populations, given specific torque model constraints?



First executable model: parsimonious ODE integration of single pulsar tracks for isolated magnetic-dipole spin-down with optional inclination decay and wind braking, embedded in a Monte Carlo population synthesis framework. Binary recycling is represented as an accretion-phase switch with an effective spin-up law. The model excludes glitches, magnetospheric state changes, variable moment of inertia, detailed companion ablation, and full binary evolution.

\subsection{Model Assumptions and Applicability}

\begin{itemize}\raggedright
\item ASM-001: Stellar geometry, electromagnetic constants, and moment of inertia are absorbed into effective braking constants K\_\allowbreak{}dipole and K\_\allowbreak{}acc. If violated: Derived constants lose direct physical interpretability and may require explicit microphysical coupling.
\item ASM-002: Magnetic inclination evolves smoothly via exponential decay without discontinuous glitches or magnetospheric state transitions. If violated: Smooth ODE integration fails to capture abrupt angular momentum transfers, invalidating continuous derivative assumptions.
\item ASM-003: Accretion torque is active exclusively during t < t\_\allowbreak{}acc with nonzero Mdot; isolated dipole torque applies otherwise. If violated: Simultaneous application of opposing torques produces unphysical period trajectories outside the validated regime.
\end{itemize}

\subsection{Symbols and Units}

\begin{center}
\small
\begin{tabular}{|p{0.12\linewidth}|p{0.15\linewidth}|p{0.46\linewidth}|p{0.17\linewidth}|}
\hline
ID & Symbol & Meaning & Unit \\ \hline
S-P & $P$ & Spin period & $\mathrm{s}$ \\ \hline
S-alpha & $\alpha$ & Magnetic inclination angle & $\mathrm{rad}$ \\ \hline
S-t & $t$ & Integration time & $\mathrm{s}$ \\ \hline
S-P0 & $P_0$ & Initial spin period & $\mathrm{s}$ \\ \hline
S-alpha0 & $\alpha_0$ & Initial inclination angle & $\mathrm{rad}$ \\ \hline
S-B0 & $B_0$ & Surface magnetic field strength & $\mathrm{T}$ \\ \hline
S-Kdip & $K_{\mathrm{dipole}}$ & Vacuum dipole braking constant & $\mathrm{s}\,\mathrm{T}^{-2}$ \\ \hline
S-eta & $\eta$ & Wind braking efficiency & $\mathrm{dimensionless}$ \\ \hline
S-tau & $\tau_\alpha$ & Inclination decay timescale & $\mathrm{s}$ \\ \hline
S-Mdot & $\dot{M}$ & Mass accretion rate & $\mathrm{kg}\,\mathrm{s}^{-1}$ \\ \hline
S-Kacc & $K_{\mathrm{acc}}$ & Accretion spin-up constant & $\mathrm{s}^{2}\,\mathrm{kg}^{-1}$ \\ \hline
S-tacc & $t_{\mathrm{acc}}$ & Accretion phase duration & $\mathrm{s}$ \\ \hline
S-Tmax & $T_{\mathrm{max}}$ & Maximum simulated age & $\mathrm{s}$ \\ \hline
S-Pdotmin & $P_{\mathrm{dot,min}}$ & Minimum detectable period derivative & $\mathrm{dimensionless}$ \\ \hline
S-Nmc & $N_{\mathrm{mc}}$ & Monte Carlo sample count & $\mathrm{dimensionless}$ \\ \hline
\end{tabular}
\end{center}

\subsection{Mechanistic Equations}

\begin{equation}\label{eq:Q1-EQ-001}
\frac{dP}{dt} = \begin{cases} -\frac{K_{\mathrm{acc}} \dot{M}}{\max(P, 0.01)} & \mathrm{if} t < t_{\mathrm{acc}} \land \dot{M} > 0 \\ +\frac{K_{\mathrm{dipole}} B_0^2 \sin^2(\alpha)(1+\eta)}{\max(P, 0.01)} & \mathrm{otherwise} \end{cases}
\end{equation}
\noindent\textbf{Q1-EQ-001 where} S-P: Spin period; S-t: Integration time; S-Kacc: Accretion spin-up constant; S-Mdot: Mass accretion rate; S-tacc: Accretion phase duration; S-Kdip: Vacuum dipole braking constant; S-B0: Surface magnetic field strength; S-alpha: Magnetic inclination angle; S-eta: Wind braking efficiency.

\begin{equation}\label{eq:Q1-EQ-002}
\frac{d\alpha}{dt} = -\frac{\alpha}{\tau_\alpha}
\end{equation}
\noindent\textbf{Q1-EQ-002 where} S-alpha: Magnetic inclination angle; S-t: Integration time; S-tau: Inclination decay timescale.

\begin{equation}\label{eq:Q1-EQ-003}
n(t) = \frac{P(t) \cdot \frac{d^2P}{dt^2}}{\left(\frac{dP}{dt}\right)^2}
\end{equation}
\noindent\textbf{Q1-EQ-003 where} S-P: Spin period.

\begin{equation}\label{eq:Q1-EQ-004}
\tau_c(t) = \frac{P(t)}{2 \cdot \frac{dP}{dt}}
\end{equation}
\noindent\textbf{Q1-EQ-004 where} S-P: Spin period.

\subsection{Initial Conditions, Boundary Conditions, and Constraints}

\paragraph{Initial Conditions}
\begin{itemize}\raggedright
\item \{'state': 'S-P', 'value': 0.5\}
\item \{'state': 'S-alpha', 'value': 0.7853981633974483\}
\end{itemize}

\paragraph{Boundary Conditions}
\begin{itemize}\raggedright
\item \{'condition': 'Integration terminates when t reaches T\_\allowbreak{}max or |dP/\allowbreak{}dt| falls below Pdot\_\allowbreak{}min', 'type': 'TERMINATION'\}
\end{itemize}

\paragraph{Objective and Constraints}
\begin{itemize}\raggedright
\item \{'objective': 'Generate distinct P-Pdot-n trajectories for isolated and recycled populations to enable statistical discrimination.'\}
\item \{'constraint': 'Sign convention: positive dP/\allowbreak{}dt denotes isolated spin-down; negative dP/\allowbreak{}dt denotes accretion spin-up. Torque branches must not activate simultaneously.'\}
\end{itemize}

\subsection{Algorithm}

\noindent\textbf{Algorithm}
\paragraph{Input}
\begin{itemize}\raggedright
\item Initial conditions (P0, alpha0)
\item Approved parameter set
\item Scenario configuration
\end{itemize}
\paragraph{Output}
\begin{itemize}\raggedright
\item Time series arrays for P(t), alpha(t)
\item Derived diagnostics n(t), tau\_\allowbreak{}c(t)
\item Population histograms
\end{itemize}
\paragraph{Steps}
\begin{itemize}\raggedright
\item Initialize state vector with P0 and alpha0
\item Select scenario parameter overrides
\item Evaluate derivative AST at current time t
\item Advance state using fixed-step integrator
\item Compute n(t) and tau\_\allowbreak{}c(t) from derivatives
\item Check termination criteria against T\_\allowbreak{}max and Pdot\_\allowbreak{}min
\item Repeat until termination or N\_\allowbreak{}mc tracks completed
\end{itemize}

\subsection{Parameters, Scenarios, and Numerical Settings}

\paragraph{Parameters}
\begin{itemize}\raggedright
\item \{'symbol\_\allowbreak{}id': 'S-P0', 'value': 0.5\}
\item \{'symbol\_\allowbreak{}id': 'S-alpha0', 'value': 0.7853981633974483\}
\item \{'symbol\_\allowbreak{}id': 'S-B0', 'value': 100000000.0\}
\item \{'symbol\_\allowbreak{}id': 'S-Kdip', 'value': 9.765625e-32\}
\item \{'symbol\_\allowbreak{}id': 'S-eta', 'value': 0.1\}
\item \{'symbol\_\allowbreak{}id': 'S-tau', 'value': 315576000000000.0\}
\item \{'symbol\_\allowbreak{}id': 'S-Mdot', 'value': 630000000000000.0\}
\item \{'symbol\_\allowbreak{}id': 'S-Kacc', 'value': 6.3e-31\}
\item \{'symbol\_\allowbreak{}id': 'S-tacc', 'value': 315576000000000.0\}
\item \{'symbol\_\allowbreak{}id': 'S-Tmax', 'value': 3155760000000000.0\}
\item \{'symbol\_\allowbreak{}id': 'S-Pdotmin', 'value': 1e-22\}
\item \{'symbol\_\allowbreak{}id': 'S-Nmc', 'value': 5000.0\}
\end{itemize}

\paragraph{Scenarios}
\begin{itemize}\raggedright
\item \{'scenario\_\allowbreak{}id': 'isolated\_\allowbreak{}dipole', 'parameter\_\allowbreak{}overrides': \{'K\_\allowbreak{}acc': 0.0, 'Mdot': 0.0, 't\_\allowbreak{}acc': 0.0\}\}
\item \{'scenario\_\allowbreak{}id': 'binary\_\allowbreak{}recycling', 'parameter\_\allowbreak{}overrides': \{'K\_\allowbreak{}acc': 5e-31\}\}
\end{itemize}

\noindent\textbf{Solver family.} Explicit Runge-Kutta 4th Order.

\subsection{Parameter Provenance and Applicability}

\noindent\textbf{Approved parameter-set identity.} 0f0174315925c8e5\allowbreak{}b143507cd77c961e\allowbreak{}1becdad36afebaa3\allowbreak{}1e7f6df12d27946c.
\begin{itemize}\raggedright
\item initial\_\allowbreak{}period (P0) \ensuremath{=} 0.5 s; role BOUNDARY\_\allowbreak{}CONDITION; status APPROVED\_\allowbreak{}USER\_\allowbreak{}INPUT; source: Birth and Evolution of Isolated Radio Pulsars; DOI 10.1086/\allowbreak{}501516; document Q1-FGK2006; locator: preprint; applicability: required\_\allowbreak{}condition\_\allowbreak{}1\ensuremath{=}For isolated or recycled neutron stars at the beginning of the simulated spin-down phase, with period measured in seconds and no glitch-induced period jump at the initial epoch.. uncertainty: not reported.
\item initial\_\allowbreak{}inclination (alpha0) \ensuremath{=} 0.7853981633974483 rad; role BOUNDARY\_\allowbreak{}CONDITION; status APPROVED\_\allowbreak{}MODEL\_\allowbreak{}ASSUMPTION; source: MODEL ASSUMPTION (not a measured or literature value); applicability: required\_\allowbreak{}condition\_\allowbreak{}1\ensuremath{=}For dipole torque models with approximately aligned magnetic and rotational axes represented by an angle in radians; treated as the starting value before inclination decay.. uncertainty: not reported.
\item initial\_\allowbreak{}surface\_\allowbreak{}magnetic\_\allowbreak{}field (B0) \ensuremath{=} 100000000.0 T; role BOUNDARY\_\allowbreak{}CONDITION; status APPROVED\_\allowbreak{}USER\_\allowbreak{}INPUT; source: Binary and Millisecond Pulsars; DOI 10.12942/\allowbreak{}lrr-2008-8; document Q1-LORIMER2008; locator: review\_\allowbreak{}article; applicability: required\_\allowbreak{}condition\_\allowbreak{}1\ensuremath{=}For isolated neutron stars with approximately dipolar surface fields inferred from timing data, expressed in tesla; not intended for magnetar-like or strongly multipolar fields.. uncertainty: not reported.
\item vacuum\_\allowbreak{}dipole\_\allowbreak{}braking\_\allowbreak{}constant (K\_\allowbreak{}dipole) \ensuremath{=} 9.765625e-32 s\_\allowbreak{}T\textasciicircum{}-2; role MODEL\_\allowbreak{}ASSUMPTION; status APPROVED\_\allowbreak{}MODEL\_\allowbreak{}ASSUMPTION; source: MODEL ASSUMPTION (not a measured or literature value); applicability: required\_\allowbreak{}condition\_\allowbreak{}1\ensuremath{=}For the chosen effective spin-down law with magnetic field in tesla and period in seconds, assuming fixed stellar geometry and no unresolved internal braking physics.. uncertainty: not reported.
\item wind\_\allowbreak{}braking\_\allowbreak{}efficiency (eta) \ensuremath{=} 0.1 dimensionless; role MODEL\_\allowbreak{}ASSUMPTION; status APPROVED\_\allowbreak{}MODEL\_\allowbreak{}ASSUMPTION; source: MODEL ASSUMPTION (not a measured or literature value); applicability: required\_\allowbreak{}condition\_\allowbreak{}1\ensuremath{=}For pulsars with magnetospheric wind torque included as a multiplicative enhancement of the dipole term; scenario switches may disable the wind contribution.. uncertainty: not reported.
\item inclination\_\allowbreak{}decay\_\allowbreak{}timescale (tau\_\allowbreak{}alpha) \ensuremath{=} 315576000000000.0 s; role MODEL\_\allowbreak{}ASSUMPTION; status APPROVED\_\allowbreak{}USER\_\allowbreak{}INPUT; source: Binary and Millisecond Pulsars; DOI 10.12942/\allowbreak{}lrr-2008-8; document Q1-LORIMER2008; locator: review\_\allowbreak{}article; applicability: required\_\allowbreak{}condition\_\allowbreak{}1\ensuremath{=}For isolated pulsars modeled with smooth inclination decay, expressed in seconds, excluding glitches and abrupt magnetospheric state changes.. uncertainty: not reported.
\item accretion\_\allowbreak{}rate (Mdot) \ensuremath{=} 630000000000000.0 kg\_\allowbreak{}s\textasciicircum{}-1; role SCENARIO\_\allowbreak{}INPUT; status APPROVED\_\allowbreak{}MODEL\_\allowbreak{}ASSUMPTION; source: MODEL ASSUMPTION (not a measured or literature value); applicability: required\_\allowbreak{}condition\_\allowbreak{}1\ensuremath{=}For recycled pulsars in a steady disk-accretion phase, expressed in kilograms per second, ignoring transient outbursts and companion ablation.. uncertainty: not reported.
\item accretion\_\allowbreak{}braking\_\allowbreak{}constant (K\_\allowbreak{}acc) \ensuremath{=} 6.3e-31 s\textasciicircum{}2\_\allowbreak{}kg\textasciicircum{}-1; role SCENARIO\_\allowbreak{}INPUT; status APPROVED\_\allowbreak{}MODEL\_\allowbreak{}ASSUMPTION; source: MODEL ASSUMPTION (not a measured or literature value); applicability: required\_\allowbreak{}condition\_\allowbreak{}1\ensuremath{=}For the effective accretion law using accretion rate in kilograms per second and dimensionless period derivative; active only during the accretion phase.. uncertainty: not reported.
\item accretion\_\allowbreak{}duration (t\_\allowbreak{}acc) \ensuremath{=} 315576000000000.0 s; role SCENARIO\_\allowbreak{}INPUT; status APPROVED\_\allowbreak{}MODEL\_\allowbreak{}ASSUMPTION; source: MODEL ASSUMPTION (not a measured or literature value); applicability: required\_\allowbreak{}condition\_\allowbreak{}1\ensuremath{=}For binary recycling tracks represented by a finite accretion interval followed by isolated spin-down, expressed in seconds.. uncertainty: not reported.
\item maximum\_\allowbreak{}simulated\_\allowbreak{}age (T\_\allowbreak{}max) \ensuremath{=} 3155760000000000.0 s; role BOUNDARY\_\allowbreak{}CONDITION; status APPROVED\_\allowbreak{}USER\_\allowbreak{}INPUT; source: Birth and Evolution of Isolated Radio Pulsars; DOI 10.1086/\allowbreak{}501516; document Q1-FGK2006; locator: preprint; applicability: required\_\allowbreak{}condition\_\allowbreak{}1\ensuremath{=}Numerical endpoint for population synthesis tracks; not a physical observable unless tied to survey selection or evolutionary assumptions.. uncertainty: not reported.
\item minimum\_\allowbreak{}detectable\_\allowbreak{}period\_\allowbreak{}derivative (Pdot\_\allowbreak{}min) \ensuremath{=} 1e-22 dimensionless; role BOUNDARY\_\allowbreak{}CONDITION; status APPROVED\_\allowbreak{}MODEL\_\allowbreak{}ASSUMPTION; source: MODEL ASSUMPTION (not a measured or literature value); applicability: required\_\allowbreak{}condition\_\allowbreak{}1\ensuremath{=}Survey-dependent detection boundary for timing observations; used as a stopping condition rather than a physical spin-down parameter.. uncertainty: not reported.
\item monte\_\allowbreak{}carlo\_\allowbreak{}sample\_\allowbreak{}count (N\_\allowbreak{}mc) \ensuremath{=} 5000.0 dimensionless; role MODEL\_\allowbreak{}ASSUMPTION; status APPROVED\_\allowbreak{}MODEL\_\allowbreak{}ASSUMPTION; source: MODEL ASSUMPTION (not a measured or literature value); applicability: required\_\allowbreak{}condition\_\allowbreak{}1\ensuremath{=}Numerical control for population synthesis; chosen for computational budget and statistical resolution rather than physical measurement.. uncertainty: not reported.
\end{itemize}

\subsection{Numerical Validation and Model-Internal Results}

\paragraph{Validation Plan}
\begin{itemize}\raggedright
\item \{'method': 'Analytical limit check: verify dalpha/\allowbreak{}dt converges to zero as t approaches infinity', 'metric': 'Relative error < 1e-6'\}
\item \{'method': 'Scenario isolation test: confirm isolated\_\allowbreak{}dipole yields strictly positive dP/\allowbreak{}dt throughout', 'metric': 'Sign consistency across all t'\}
\item \{'method': 'Recycling transition test: verify dP/\allowbreak{}dt switches sign at t\_\allowbreak{}acc and returns to positive post-accretion', 'metric': 'Branch activation count matches scenario definition'\}
\end{itemize}

\begin{itemize}
\item Execution sim-ef92f26db7c14196aaab75759811335e; NUMERICAL\_SIMULATION; SIMULATED; NOT\_EMPIRICAL; relation CONSTRAINED; summary: Within the audited reduced ODE, the isolated track increases from P\ensuremath{=}0.500 s to 0.662 s, whereas the recycling track reaches P\_min\ensuremath{=}0.226 s during accretion and ends at 0.281 s after resumed dipole spin-down. This establishes model-internal directional separation, but does not by itself validate population-level n-tau\_c discrimination because the executable MathIR contains two deterministic tracks rather than the documented Monte Carlo outer loop.
\end{itemize}

\subsection{Hypothesis Iteration Lineage}

\begin{itemize}\raggedright
\item v0: CONSTRAINED (Within the audited reduced ODE, the isolated track increases from P\ensuremath{=}0.500 s to 0.662 s, whereas the recycling track reaches P\_\allowbreak{}min\ensuremath{=}0.226 s during accretion and ends at 0.281 s after resumed dipole spin-down. This establishes model-internal directional separation, but does not by itself validate population-level n-tau\_\allowbreak{}c discrimination because the executable MathIR contains two deterministic tracks rather than the documented Monte Carlo outer loop.)
\end{itemize}

\subsection{Limitations, Non-Empirical Disclosure, and References}

\noindent All values in this section are model-internal numerical simulations (NUMERICAL\_SIMULATION; SIMULATED; NOT\_EMPIRICAL), not empirical observations.

\paragraph{Limitations}
\begin{itemize}\raggedright
\item Excludes glitch events, magnetospheric state transitions, and variable moment of inertia.
\item Assumes fixed stellar geometry and absorbs unresolved electromagnetic factors into effective constants.
\item Observational selection effects and survey detection thresholds are not dynamically coupled to the simulation loop.
\end{itemize}

\paragraph{References}
\begin{itemize}\raggedright
\item \{'citation': 'Feng, G., et al. (2006). Birth and Evolution of Isolated Radio Pulsars. ApJ, 646, 381. DOI: 10.1086/\allowbreak{}501516', 'source': 'Q1-FGK2006'\}
\item \{'citation': 'Lorimer, D. R. (2008). Binary and Millisecond Pulsars. Living Rev. Relativity, 11, 8. DOI: 10.12942/\allowbreak{}lrr-2008-8', 'source': 'Q1-LORIMER2008'\}
\end{itemize}

\section{Q2 (final: v1)}

\noindent\textbf{Abstract—} A spatially resolved spherical radial thermal model for a massive neutron star core is specified as an executable PDE branch. The model replaces the parent lumped thermal state with a one-dimensional radial heat-transport field in the canonical coordinate x, where x represents the physical radial coordinate r normalized by the stellar radius. Phase-transition latent heat and modified neutrino cooling are represented through parameterized volumetric source terms. Scenario A is standard cooling without latent heat, Scenario B adds localized latent-heat release, and Scenario C modifies the cooling coefficient to emulate superfluid suppression. The declared observable outputs are radial temperature trajectories and a surface flux proxy at x\ensuremath{=}1.

\subsection{Question and Scope}

\noindent\textbf{Question.} How does the radial thermal evolution of a massive neutron star differ when spin-down-associated core compression triggers a localized phase-transition heat source, compared with standard cooling and with a superfluid-modified cooling scenario?



Executable spherical radial thermal PDE branch for controlled numerical simulation. The model covers radial temperature evolution, a parameterized phase-transition heat source, volumetric cooling, center regularity, and a fixed scaled surface temperature boundary. It does not solve spin dynamics, density evolution, full microphysics, atmosphere radiative transfer, or observational inference.

\subsection{Model Assumptions and Applicability}

\begin{itemize}\raggedright
\item Q2-ASM-001: The star is spherically symmetric and the executable PDE uses the canonical coordinate x as the physical radial coordinate r normalized by the stellar radius. If violated: Non-spherical structure or multidimensional transport would not be represented, biasing radial temperature and surface-flux estimates.
\item Q2-ASM-002: Spin-down induced core compression and phase-transition onset are represented by a prescribed localized latent-heat source rather than dynamically solved spin, density, or moment-of-inertia equations. If violated: Feedback between rotation, compression, and transition timing is omitted, so scenario attribution becomes phenomenological.
\item Q2-ASM-003: Thermal microphysics is reduced to scaled constant transport, heat-capacity, and source coefficients. If violated: Temperature- and density-dependent conductivity, heat capacity, and neutrino emissivity would change cooling rates and scenario separation.
\item Q2-ASM-004: The surface boundary is approximated by a fixed scaled surface temperature representing atmosphere and radiative-transfer effects. If violated: An inaccurate surface boundary would alter the surface gradient, flux proxy, and inferred luminosity.
\item Q2-ASM-005: Units are normalized so that the temperature field is scaled by a reference core temperature, x\ensuremath{=}1 corresponds to the stellar radius, and t\ensuremath{=}1 corresponds to a characteristic cooling time. If violated: Direct physical-unit comparison requires rescaling; mismatched scaling would distort comparison with cooling curves.
\end{itemize}

\subsection{PDE Model Class}

\noindent\textbf{System type.} SPHERICAL\_RADIAL\_THERMAL.

\noindent\textbf{Trusted adapter.} spherical\_radial\_thermal.

\subsection{Spatial Domain and Field Variables}

\paragraph{Domain and mesh}
\begin{itemize}\raggedright
\item Domain: x\ensuremath{=}[0.0, 1.0]
\item Grid: nx\ensuremath{=}33
\end{itemize}

\paragraph{Fields}
\begin{itemize}\raggedright
\item T (T); unit 10\textasciicircum{}9 K
\end{itemize}

\subsection{PDE Initial and Boundary Conditions}

\paragraph{Initial condition}
\begin{itemize}\raggedright
\item Sampled values on the normalized mesh.
\item Time span: 0.0; 10.0
\end{itemize}

\paragraph{Boundary condition types}
\begin{itemize}\raggedright
\item left: SPHERICAL\_\allowbreak{}ORIGIN\_\allowbreak{}REGULARITY
\item right: DIRICHLET
\end{itemize}

\subsection{Discretization and Stability}

\noindent\textbf{Discretization.} FINITE\_DIFFERENCE\_SPHERICAL\_RADIAL.

\noindent\textbf{Time integration.} EXPLICIT\_EULER; time step 0.005.

\subsection{Numerical Verification}

\noindent The fixed PDE adapter reports finite-field checks and the family-specific stability or linear-residual diagnostic. A result is qualified only when the declared checks pass; all outcomes remain model-internal simulations.

\subsection{Symbols and Units}

\begin{center}
\small
\begin{tabular}{|p{0.12\linewidth}|p{0.15\linewidth}|p{0.46\linewidth}|p{0.17\linewidth}|}
\hline
ID & Symbol & Meaning & Unit \\ \hline
T & $T$ & Radial temperature field scaled by a reference core temperature & $10^9 K$ \\ \hline
x & $x$ & Canonical radial coordinate representing physical radius r normalized by stellar radius & $\mathrm{R}\,\mathrm{NS}$ \\ \hline
t & $t$ & Time scaled by a characteristic cooling time & $\mathrm{t}\,\mathrm{char}$ \\ \hline
source & $S$ & Volumetric heat source and sink term in scaled temperature per time & $10^9 K/t_char$ \\ \hline
kappa & $\kappa$ & Effective radial thermal diffusion coefficient in scaled units & $R_NS^2/t_char$ \\ \hline
c\_v & $c_v$ & Volumetric heat capacity in scaled units & energy/(R\_NS\textasciicircum{}3*10\textasciicircum{}9 K) \\ \hline
latent\_amp & $L_{\mathrm{amp}}$ & Amplitude of latent-heat release from the phase-transition layer & $10^9 K/t_char$ \\ \hline
cooling\_amp & $\epsilon_{\mathrm{amp}}$ & Amplitude of volumetric neutrino cooling coefficient & $1/t_char$ \\ \hline
transition\_center & $x_{\mathrm{tr}}$ & Normalized radial center of the phase-transition layer & $\mathrm{R}\,\mathrm{NS}$ \\ \hline
transition\_width & $\sigma_{\mathrm{tr}}$ & Normalized radial width of the phase-transition layer & $\mathrm{R}\,\mathrm{NS}$ \\ \hline
decay\_rate & $\lambda$ & Inverse time scale for transition-source decay & $1/t_char$ \\ \hline
surface\_temperature & $T_s$ & Scaled surface temperature imposed at x\ensuremath{=}1 & $10^9 K$ \\ \hline
\end{tabular}
\end{center}

\subsection{Mechanistic Equations}

\begin{equation}\label{eq:Q2-EQ-001}
c_v \frac{\partial T}{\partial t} = \frac{1}{x^2}\frac{\partial}{\partial x}\left(x^2 \kappa \frac{\partial T}{\partial x}\right) + S
\end{equation}
\noindent\textbf{Q2-EQ-001 where} c\_v: Volumetric heat capacity in scaled units; T: Radial temperature field scaled by a reference core temperature; t: Time scaled by a characteristic cooling time; x: Canonical radial coordinate representing physical radius r normalized by stellar radius; kappa: Effective radial thermal diffusion coefficient in scaled units; source: Volumetric heat source and sink term in scaled temperature per time.

\begin{equation}\label{eq:Q2-EQ-002}
S = L_{\mathrm{amp}}\exp\left[-\frac{(x-x_{\mathrm{tr}})^2}{2\sigma_{\mathrm{tr}}^2}\right]\exp(-\lambda t) - \epsilon_{\mathrm{amp}} T
\end{equation}
\noindent\textbf{Q2-EQ-002 where} source: Volumetric heat source and sink term in scaled temperature per time; latent\_amp: Amplitude of latent-heat release from the phase-transition layer; x: Canonical radial coordinate representing physical radius r normalized by stellar radius; transition\_center: Normalized radial center of the phase-transition layer; transition\_width: Normalized radial width of the phase-transition layer; decay\_rate: Inverse time scale for transition-source decay; t: Time scaled by a characteristic cooling time; cooling\_amp: Amplitude of volumetric neutrino cooling coefficient; T: Radial temperature field scaled by a reference core temperature.

\begin{equation}\label{eq:Q2-EQ-003}
\left.\frac{\partial T}{\partial x}\right|_{x=0}=0,  T(1,t)=T_s
\end{equation}
\noindent\textbf{Q2-EQ-003 where} T: Radial temperature field scaled by a reference core temperature; x: Canonical radial coordinate representing physical radius r normalized by stellar radius; t: Time scaled by a characteristic cooling time; surface\_temperature: Scaled surface temperature imposed at x\ensuremath{=}1.

\subsection{Initial Conditions, Boundary Conditions, and Constraints}

\paragraph{Initial Conditions}
\begin{itemize}\raggedright
\item \{'condition\_\allowbreak{}id': 'Q2-IC-001', 'field\_\allowbreak{}id': 'T', 'profile': 'UNIFORM', 'description': 'Uniform scaled initial core temperature across the radial domain.'\}
\end{itemize}

\paragraph{Boundary Conditions}
\begin{itemize}\raggedright
\item \{'boundary\_\allowbreak{}id': 'Q2-BC-001', 'side': 'left', 'type': 'SPHERICAL\_\allowbreak{}ORIGIN\_\allowbreak{}REGULARITY', 'description': 'Regularity at the spherical origin x\ensuremath{=}0.'\}
\item \{'boundary\_\allowbreak{}id': 'Q2-BC-002', 'side': 'right', 'type': 'DIRICHLET', 'value\_\allowbreak{}parameter': 'surface\_\allowbreak{}temperature', 'description': 'Fixed scaled surface temperature at x\ensuremath{=}1.'\}
\end{itemize}

\paragraph{Objective and Constraints}
\begin{itemize}\raggedright
\item \{'objective\_\allowbreak{}id': 'Q2-OBJ-001', 'statement': 'Compare radial temperature trajectories and the surface flux proxy among standard, phase-transition, and superfluid-modified scenarios.'\}
\item \{'objective\_\allowbreak{}id': 'Q2-OBJ-002', 'statement': 'Determine whether localized latent-heat release produces a distinguishable delay or deviation in cooling relative to standard cooling.'\}
\item \{'constraint\_\allowbreak{}id': 'Q2-CON-001', 'statement': 'Use only the registered spherical radial thermal PDE adapter with a uniform grid and explicit time integration.'\}
\item \{'constraint\_\allowbreak{}id': 'Q2-CON-002', 'statement': 'Maintain non-negative diffusion coefficient, positive heat capacity, and finite field values throughout the declared time span.'\}
\end{itemize}

\subsection{Algorithm}

\noindent\textbf{Algorithm}
\paragraph{Input}
\begin{itemize}\raggedright
\item scenario parameter values
\item spatial\_\allowbreak{}domain
\item grid
\item initial\_\allowbreak{}condition
\item boundary\_\allowbreak{}conditions
\end{itemize}
\paragraph{Output}
\begin{itemize}\raggedright
\item radial temperature field T(x,t)
\item surface flux proxy at x\ensuremath{=}1
\item scenario comparison diagnostics
\end{itemize}
\paragraph{Steps}
\begin{itemize}\raggedright
\item Select scenario parameter values.
\item Initialize the uniform scaled temperature field.
\item Apply spherical origin regularity and the surface Dirichlet condition.
\item Advance the spherical radial heat equation with the registered explicit integrator.
\item Evaluate stability and finiteness checks.
\item Extract radial profiles and the surface flux proxy.
\item Compare scenario tracks against the declared falsification criteria.
\end{itemize}

\subsection{Parameters, Scenarios, and Numerical Settings}

\paragraph{Parameters}
\begin{itemize}\raggedright
\item \{'parameter\_\allowbreak{}id': 'Q2-PRM-001', 'symbol\_\allowbreak{}id': 'kappa', 'base\_\allowbreak{}value': 0.01, 'evidence\_\allowbreak{}status': 'model\_\allowbreak{}assumption', 'source': 'Scaled effective radial thermal diffusion chosen for explicit spherical stability.'\}
\item \{'parameter\_\allowbreak{}id': 'Q2-PRM-002', 'symbol\_\allowbreak{}id': 'c\_\allowbreak{}v', 'base\_\allowbreak{}value': 1.0, 'evidence\_\allowbreak{}status': 'model\_\allowbreak{}assumption', 'source': 'Scaled volumetric heat capacity reference.'\}
\item \{'parameter\_\allowbreak{}id': 'Q2-PRM-003', 'symbol\_\allowbreak{}id': 'latent\_\allowbreak{}amp', 'base\_\allowbreak{}value': 0.0, 'evidence\_\allowbreak{}status': 'scenario\_\allowbreak{}input', 'source': 'No latent heat in the standard scenario; positive in phase-transition scenarios.'\}
\item \{'parameter\_\allowbreak{}id': 'Q2-PRM-004', 'symbol\_\allowbreak{}id': 'cooling\_\allowbreak{}amp', 'base\_\allowbreak{}value': 0.02, 'evidence\_\allowbreak{}status': 'model\_\allowbreak{}assumption', 'source': 'Scaled neutrino cooling coefficient.'\}
\item \{'parameter\_\allowbreak{}id': 'Q2-PRM-005', 'symbol\_\allowbreak{}id': 'transition\_\allowbreak{}center', 'base\_\allowbreak{}value': 0.35, 'evidence\_\allowbreak{}status': 'model\_\allowbreak{}assumption', 'source': 'Normalized core radius where the phase transition is concentrated.'\}
\item \{'parameter\_\allowbreak{}id': 'Q2-PRM-006', 'symbol\_\allowbreak{}id': 'transition\_\allowbreak{}width', 'base\_\allowbreak{}value': 0.12, 'evidence\_\allowbreak{}status': 'model\_\allowbreak{}assumption', 'source': 'Normalized radial width of the transition layer.'\}
\item \{'parameter\_\allowbreak{}id': 'Q2-PRM-007', 'symbol\_\allowbreak{}id': 'decay\_\allowbreak{}rate', 'base\_\allowbreak{}value': 0.6, 'evidence\_\allowbreak{}status': 'model\_\allowbreak{}assumption', 'source': 'Inverse time scale for latent-heat release decay.'\}
\item \{'parameter\_\allowbreak{}id': 'Q2-PRM-008', 'symbol\_\allowbreak{}id': 'surface\_\allowbreak{}temperature', 'base\_\allowbreak{}value': 0.1, 'evidence\_\allowbreak{}status': 'model\_\allowbreak{}assumption', 'source': 'Scaled surface temperature representing the atmosphere boundary.'\}
\end{itemize}

\paragraph{Scenarios}
\begin{itemize}\raggedright
\item \{'scenario\_\allowbreak{}id': 'Q2-SCN-A', 'name': 'Standard cooling', 'description': 'No phase-transition latent heat and standard scaled neutrino cooling.', 'parameter\_\allowbreak{}overrides': \{'latent\_\allowbreak{}amp': 0.0, 'cooling\_\allowbreak{}amp': 0.02, 'surface\_\allowbreak{}temperature': 0.1\}\}
\item \{'scenario\_\allowbreak{}id': 'Q2-SCN-B', 'name': 'Phase transition cooling', 'description': 'Localized latent-heat release associated with spin-down induced core compression.', 'parameter\_\allowbreak{}overrides': \{'latent\_\allowbreak{}amp': 0.3, 'cooling\_\allowbreak{}amp': 0.02, 'surface\_\allowbreak{}temperature': 0.1\}\}
\item \{'scenario\_\allowbreak{}id': 'Q2-SCN-C', 'name': 'Phase transition with superfluidity', 'description': 'Latent-heat release with suppressed neutrino cooling representing superfluid modification.', 'parameter\_\allowbreak{}overrides': \{'latent\_\allowbreak{}amp': 0.3, 'cooling\_\allowbreak{}amp': 0.01, 'surface\_\allowbreak{}temperature': 0.1\}\}
\end{itemize}

\noindent\textbf{Solver family.} spherical\_radial\_thermal.

\subsection{Parameter Provenance and Applicability}

\noindent\textbf{Parameter status.} This legacy model used inline assumptions; it has no approved external parameter set.

\subsection{Numerical Validation and Model-Internal Results}

\paragraph{Validation Plan}
\begin{itemize}\raggedright
\item \{'check\_\allowbreak{}id': 'Q2-VAL-001', 'description': 'Preflight stability: the explicit spherical diffusion number remains below the registered limit.', 'status': 'required\_\allowbreak{}before\_\allowbreak{}execution'\}
\item \{'check\_\allowbreak{}id': 'Q2-VAL-002', 'description': 'Scenario separation: phase-transition scenarios produce different radial temperature histories and surface flux proxy values from the standard scenario.', 'status': 'required\_\allowbreak{}after\_\allowbreak{}execution'\}
\item \{'check\_\allowbreak{}id': 'Q2-VAL-003', 'description': 'Boundary consistency: center regularity and surface Dirichlet condition are satisfied without non-finite values.', 'status': 'required\_\allowbreak{}after\_\allowbreak{}execution'\}
\item \{'check\_\allowbreak{}id': 'Q2-VAL-004', 'description': 'Observational relevance: the surface flux proxy is mapped to luminosity only after external scaling and calibration.', 'status': 'postprocessing'\}
\end{itemize}

\begin{itemize}
\item Execution sim-1b3a94f9eebb4426b07d75c423ce2def; NUMERICAL\_SIMULATION; SIMULATED; NOT\_EMPIRICAL; relation INCONCLUSIVE; summary: Q2 v1 baseline spherical-radial thermal PDE completed with finite output, explicit stability, spherical-origin regularity, and boundary checks passing; this baseline-only run does not establish scenario separation or empirical pulsar validation.
\end{itemize}

\subsection{Hypothesis Iteration Lineage}

\begin{itemize}\raggedright
\item v0: SUPPORTED\_\allowbreak{}WITHIN\_\allowbreak{}MODEL (The audited reduced ODE crosses rho\_\allowbreak{}crit\ensuremath{=}5.12e17 kg/\allowbreak{}m\textasciicircum{}3 at approximately 9.85e11 s in all spin-down cases. At 3*tau\_\allowbreak{}cool, the standard, phase-transition, and phase-transition-plus-superfluid scenarios yield surface temperatures of about 4.98e4 K, 5.32e4 K, and 2.36e5 K, respectively; phase\_\allowbreak{}fraction reaches 1 only in the two transition scenarios. The mechanism therefore produces distinguishable model-internal thermal tracks, conditional on the effective cooling, latent-heat, and compression proxies.)
\item v1: INCONCLUSIVE (Q2 v1 baseline spherical-radial thermal PDE completed with finite output, explicit stability, spherical-origin regularity, and boundary checks passing; this baseline-only run does not establish scenario separation or empirical pulsar validation.)
\end{itemize}

\subsection{Limitations, Non-Empirical Disclosure, and References}

\noindent All values in this section are model-internal numerical simulations (NUMERICAL\_SIMULATION; SIMULATED; NOT\_EMPIRICAL), not empirical observations.

\paragraph{Limitations}
\begin{itemize}\raggedright
\item The model does not solve spin evolution, moment of inertia, or central density dynamically; spin-down compression is parameterized by a prescribed source.
\item Microphysical equation-of-state effects, latent heat, neutrino emissivity, and superfluid suppression are represented by scaled coefficients rather than density-temperature tables.
\item The surface boundary is simplified; atmosphere radiative transfer and gravitational redshift are not solved.
\item Spherical symmetry and a uniform radial grid exclude multidimensional effects and sharp phase-boundary dynamics.
\item No observational data are ingested; comparison to X-ray cooling curves requires external calibration and uncertainty treatment.
\end{itemize}

\paragraph{References}
\begin{itemize}\raggedright
\item Standard neutron-star cooling theory and cooling-curve literature for model comparison.
\item X-ray observations of isolated neutron stars and massive pulsars from Chandra and XMM-Newton archives as external comparator context.
\item Equation-of-state and phase-transition literature for neutron-star cores as parameter motivation.
\end{itemize}

\end{document}
